\section{Introduction}

Mammals confronted with food scarcity reduce whole-body energy expenditure, and the nature of these adaptations is known to differ between the sexes. Studies of peripheral physiology have repeatedly shown that female rodents preserve fat mass and bodyweight better than males under comparable restriction, while simultaneously suppressing energy-costly reproductive functions such as ovulation and maintenance of uterine tissue~\cite{lovdel2024hsd11b1,lopezdominguez2016omega3,greenfield2025lifespan,ryu2023timerestricted,jeromson2024gdf15}. Such sexual dimorphism is thought to reflect the differential contribution of adipose stores versus reproductive investment to fitness in males and females~\cite{nestor2014crosstalk}.

In striking contrast, far less is understood about whether and how the brain participates in these sex-specific energy-saving programs. Neural tissue is metabolically expensive: per unit mass, grey matter consumes more ATP than almost any peripheral organ, with a large fraction of that cost attributable to action potentials, postsynaptic currents, and the Na$^+$/K$^+$ pumping needed to restore resting potentials after signalling~\cite{attwell2001energy,niven2008energy}. Evolutionary analyses argue that this high metabolic price has shaped the architecture of sensory systems toward energy-efficient codes~\cite{niven2008energy}. If the cortex can flexibly lower its energy budget during periods of scarcity, this would constitute a meaningful contribution to whole-organism metabolic adaptation.

The primary visual cortex (V1) of the mouse provides a tractable testbed for this question. Previous work in male mice demonstrated that moderate food restriction reduces V1 ATP usage during visual processing and produces a measurable loss of orientation tuning precision, and that this adaptation depends on decreased circulating leptin, an adipocyte-derived hormone~\cite{zhang1994positional}. Leptin is widely regarded as a key signal by which the brain senses peripheral energy status, and it couples to downstream intracellular pathways including AMP-activated protein kinase (AMPK), mTOR signalling, and peroxisome proliferator-activated receptor alpha (PPAR$\alpha$)-regulated fatty acid oxidation~\cite{steinberg2023ampk,laplante2013mtorc1,yao2022ampk,su2019hsf1,ward2016ampk}. In male animals, the circuitry linking fat stores to cortical function is therefore reasonably well established.

Whether the corresponding loop operates in females has remained an open question, despite its obvious relevance. Food restriction protocols are routinely used in systems neuroscience to motivate animals during behavioural tasks that probe cortical function~\cite{scott2023cisplatin,heinz2021foraging}, and the resulting metabolic state is often assumed to be comparable across sexes. Yet peripheral studies suggest that females may not enter the same metabolic state as males for an equivalent relative loss of bodyweight, because their fat and leptin dynamics differ~\cite{nestor2014crosstalk,kwon2014gpr30,elsehrawy2025eralpha}. A cortical equivalent of this asymmetry, if it exists, would have direct consequences for how behavioural data from male and female animals should be interpreted, and for how dietary interventions aimed at neurological disease are designed.

Here we address this gap directly. Using adult male and female mice food-restricted to 85\% of free-feeding bodyweight over two to three weeks, we combine biochemical assays of AMPK Thr172 phosphorylation and PPAR$\alpha$ DNA-binding activity, bulk RNA sequencing of V1 tissue analysed with DESeq2~\cite{love2014moderated} and Gene Set Enrichment Analysis~\cite{subramanian2005gsea}, two-photon FRET imaging of ATP dynamics using the ATeam1.03YEMK sensor~\cite{imamura2009visualization,tsuyama2013ateam,ando2012visualization,kishikawa2011ateam}, and GCaMP6s-based calcium imaging of V1 layer 2/3 neurons during drifting-grating presentation~\cite{mcclure2019puretones,petousakis2023apical}. We complement classical frequentist tests with Bayes factor analysis~\cite{sekulovski2024bayes} to assess the strength of evidence for or against sex-specific effects. Our central finding is that V1 energy regulation and coding precision are largely preserved in females under food restriction while they are robustly reduced in males, a pattern that parallels peripheral sex differences in leptin dynamics and that extends the concept of sex-specific metabolic resilience from the periphery to the neocortex.

\section{Related Work}

\subsection{Brain energetics and its constraints on cortical function}
Quantitative energy budgets of the cortex indicate that action potentials and postsynaptic glutamate-receptor currents dominate neuronal ATP expenditure, with smaller contributions from resting potentials and transmitter recycling~\cite{attwell2001energy}. Because even small changes in firing rate translate into substantial changes in oxygen consumption, cortical circuits face strong selective pressure to balance information coding against metabolic cost~\cite{niven2008energy}. These considerations motivate the hypothesis that, when peripheral energy becomes limiting, cortical circuits should be configured to shed metabolically expensive components of their activity.

\subsection{AMPK, mTOR, and PPAR$\alpha$ as intracellular energy regulators}
AMPK is activated by rising AMP/ATP and ADP/ATP ratios and by phosphorylation of threonine 172, after which it suppresses anabolic ATP-consuming processes (including mTORC1 signalling) and promotes catabolism and oxidative phosphorylation~\cite{steinberg2023ampk,ward2016ampk,su2019hsf1,yao2022ampk}. mTORC1 itself couples growth-factor and nutrient signals to protein synthesis, lipid biogenesis, and transcriptional regulation of metabolic programs, and its down-regulation is associated with reduced neuronal biosynthetic load~\cite{laplante2013mtorc1}. PPAR$\alpha$, a nuclear receptor whose activity rises under fasting and lipolytic states, coordinates fatty-acid oxidation downstream of leptin signalling~\cite{lopezdominguez2016omega3}. Together, these three pathways form a canonical axis for sensing and responding to cellular energy stress, but most direct measurements of their regulation have been made in peripheral tissue.

\subsection{Leptin, oestrogens, and sex differences in energy homeostasis}
Leptin, first identified by positional cloning of the $ob$ gene, is secreted in proportion to fat mass and is a major peripheral signal of energy availability to the brain~\cite{zhang1994positional}. Oestrogens interact with this signal at multiple levels: they suppress food intake, enhance leptin-induced STAT3 phosphorylation in the hypothalamus through receptors including GPR30, and promote mitochondrial activity~\cite{kwon2014gpr30,nestor2014crosstalk,elsehrawy2025eralpha}. Peripheral studies in mice find that caloric restriction affects males and females differently, and that many of these differences attenuate with age as oestrogen levels decline~\cite{lovdel2024hsd11b1,lopezdominguez2016omega3,greenfield2025lifespan,jeromson2024gdf15,ryu2023timerestricted}. Whether these sex-biased peripheral programs propagate into the neocortex has not been systematically tested.

\subsection{Two-photon imaging of cortical energy use and coding}
Genetically encoded FRET-based ATP sensors (the ATeam family) have enabled direct in vivo measurement of cellular ATP dynamics in model organisms and in vertebrate brain~\cite{imamura2009visualization,tsuyama2013ateam,ando2012visualization,kishikawa2011ateam}. Combined with focal application of ATP-synthesis inhibitors, the rate at which the FRET signal decays provides a proxy for net ATP usage during a defined behavioural epoch. In parallel, calcium indicators of the GCaMP family permit measurement of orientation tuning and direction selectivity in V1 layer 2/3 neurons~\cite{petousakis2023apical,mcclure2019puretones}, enabling direct comparison of coding precision across metabolic states. Our study applies these complementary readouts to test, for the first time to our knowledge, whether their sex-specific modulation matches the sex-specific peripheral metabolic response.

\subsection{Statistical evaluation of null effects}
Conventional null-hypothesis significance testing does not distinguish between absence of evidence and evidence of absence. Bayes factor analysis addresses this by quantifying the relative support the data provide to alternative and null hypotheses~\cite{sekulovski2024bayes}. When results in one sex are significant and in the other sex are not, Bayes factors provide additional information about whether the non-significant result reflects a true null or an underpowered experiment. We therefore report both frequentist tests and Bayes factors for our key outcomes.
